\clearpage
\begin{trapbox}

\begin{minipage}{\linewidth}
\textbf{\textcolor{CTrap}{易错一：一代入不能算，就马上分别求导}}
\begin{center}\begin{tikzpicture}[font=\small,>=stealth,line join=round]
\draw[gray!45,rounded corners=3pt] (0,0) rectangle (14.2,9.4);
\node at (7.1,8.8) {不能代入不等于可以求导：分子并未趋于零};
\begin{scope}[shift={(3.5,4.4)},x=2cm,y=.6cm]
\draw[->,gray] (-1.5,0)--(1.65,0) node[above] {$t=x-1$};
\draw[->,gray] (0,-4.5)--(0,6.2) node[left] {$y$};
\draw[blue!75!black,very thick,samples=80,domain=-1.3:-.22] plot (\x,{1+1/\x});
\draw[blue!75!black,very thick,samples=80,domain=.22:1.3] plot (\x,{1+1/\x});
\draw[orange!85!black,thick,dashed] (-1.3,1)--(1.3,1);
\node[above left=8pt] at (0,1) {$1$};
\end{scope}
\node[anchor=west,align=left] at (8.8,5.2)
{蓝：$\dfrac{x}{x-1}=1+\dfrac1t$\\[18pt]
橙：误用后得到 $1$\\[18pt]
$t\to0^+:\ +\infty$\\[13pt]
$t\to0^-:\ -\infty$};
\node[text width=13.2cm,align=center] at (7.1,.65) {原式是 $1/0$ 型；分别求导得到的常数线，不能代表它的目标极限。};
\end{tikzpicture}
\end{center}
\end{minipage}



错误写法：
\[
\lim_{x\to1}\frac{x}{x-1}
\overset{\text{错}}{=}\lim_{x\to1}\frac11=1.
\]
实际当 $x\to1^+$ 时原式趋向 $+\infty$，当 $x\to1^-$ 时趋向 $-\infty$。
\textbf{改正：}先看分子、分母到底分别趋于什么；$\frac10$ 不是 $\frac00$。

\medskip
\textbf{\textcolor{CTrap}{易错二：把“分别求导”写成“整个分式求导”}}

\[
\underbrace{\frac{f'(x)}{g'(x)}}_{\text{法则里比较的式子}}
\qquad\text{与}\qquad
\underbrace{\left(\frac{f(x)}{g(x)}\right)'}_{\text{整个分式的导数}}
\]
通常\textbf{不是同一回事}。例如 $\frac{x^2}{x}=x$，整个分式的导数是 $1$，
分别求导后的比值是 $\frac{2x}{1}=2x$。

\medskip
\textbf{\textcolor{CTrap}{易错三：一次求导后不重新检查}}

例如求 $\displaystyle\lim_{x\to0}\frac{x^2}{x}$：
\[
\lim_{x\to0}\frac{x^2}{x}
=\lim_{x\to0}\frac{2x}{1}=0.
\]
此时已经算出 $0$，\textbf{不能再次硬求导}。
每次再用，都必须重新确认新式子满足条件。

\medskip
\textbf{\textcolor{CTrap}{易错四：没有确认分母的导数可以放在分母}}

若 $g'(x)=0$，就不能随意写 $\frac{f'(x)}{g'(x)}$。
例如 $g(x)=5$ 时 $g'(x)=0$，此时本来就不属于“$0$ 比 $0$”或“两边都无限增大”的情形。
\textbf{改正：}确认在目标位置附近（不包括目标位置本身）$g'(x)\ne0$，并确认求导后的比值有确定结果。

\medskip
\begin{minipage}{\linewidth}
\textbf{\textcolor{CTrap}{易错五：求导后的极限不存在，就认定原式也不存在}}
\begin{center}\begin{tikzpicture}[font=\small,>=stealth,line join=round]
\draw[gray!45,rounded corners=3pt] (0,0) rectangle (14.2,9.4);
\node at (7.1,8.8) {原比值趋于一个数，求导后的比值却可能一直振荡};
\begin{scope}[shift={(1.4,2.6)},x=.36cm,y=2.0cm]
\draw[->,gray] (0,0)--(16.5,0) node[above] {$x$};
\draw[->,gray] (0,0)--(0,2.5) node[left] {$y$};
\draw[gray,dotted] (0,1)--(16,1);\node[left=10pt] at (0,1) {$1$};
\draw[blue!75!black,thick,domain=2:16,samples=160] plot (\x,{1+sin(deg(\x))/\x});
\draw[orange!85!black,dashed,thick,domain=2:16,samples=160] plot (\x,{1+cos(deg(\x))});
\end{scope}
\node[anchor=west,align=left] at (8.6,5.1)
{蓝：$1+\dfrac{\sin x}{x}\to1$\\[18pt]
橙：$1+\cos x$ 无极限\\[18pt]
原式偏离 $1$ 的幅度\\[8pt]不超过 $1/x$。};
\node[text width=13.2cm,align=center] at (7.1,.65) {蓝线越靠右越贴近高度 $1$；橙线仍在 $0$ 与 $2$ 之间反复变化。};
\end{tikzpicture}
\end{center}
\end{minipage}



\[
\frac{x+\sin x}{x}=1+\frac{\sin x}{x}\longrightarrow1
\qquad(x\to+\infty).
\]
但分别求导后为 $1+\cos x$，它没有极限。
\textbf{改正：}这说明求导法在此不能给出答案，不是原式没有极限。原式可由 $|\sin x|\le1$ 直接判断。

\medskip
\textbf{\textcolor{CTrap}{易错六：忘了左右两边可能不同}}

\[
\frac{x-1}{(x-1)^2}=\frac1{x-1}.
\]
$x\to1^+$ 时趋向 $+\infty$，$x\to1^-$ 时趋向 $-\infty$，所以不能说双侧极限等于 $+\infty$。
题目若只问从右边或左边靠近，要按对应方向作答。

\medskip
\textbf{\textcolor{CTrap}{易错七：为了用法则，反而把简单题做复杂}}

\[
\frac{x^2-9}{x-3}=x+3\quad(x\ne3).
\]
这类题目因式分解更快。含根号先考虑有理化，含三角函数可考虑已知基本极限。
\textbf{考试提醒：}若课堂尚未正式学习洛必达法则，书面证明与考试作答应以课程允许的方法为准。

\end{trapbox}
